% TeleCon connection modes — what the phone setting means.
% Build: pdflatex -output-directory=../build telecon_connection_modes.tex
% Delivered PDF: ../build/connection_modes_eng.pdf
\documentclass[12pt,a4paper]{article}

\usepackage[margin=2.4cm]{geometry}
\usepackage[T1]{fontenc}
\usepackage[utf8]{inputenc}
\usepackage{lmodern}
\usepackage{booktabs}
\usepackage{hyperref}
\usepackage{xcolor}
\usepackage{colortbl}
\usepackage{array}
\usepackage{enumitem}
\usepackage{needspace}
\usepackage{etoolbox}

\preto\section{\Needspace{8\baselineskip}}

\hypersetup{
  colorlinks=true,
  linkcolor=black,
  urlcolor=blue,
  pdftitle={TeleCon connection modes},
  pdfauthor={TeleCon}
}

\makeatletter
\g@addto@macro\UrlBreaks{\do\_\do\-\do\/\do\:}
\makeatother
\newcommand{\code}[1]{\nolinkurl{#1}}

\definecolor{TblHead}{HTML}{0F6F73}
\definecolor{TblChan}{HTML}{D5F5E3}
\definecolor{TblTelem}{HTML}{D4E6F1}
\definecolor{TblKnob}{HTML}{F5CBA7}

\newcommand{\thc}[1]{\textcolor{white}{\bfseries #1}}
\newcommand{\chip}[2]{\colorbox{#1}{\strut\,\scriptsize #2\,}}
\newcommand{\ui}[1]{\textsf{\textbf{#1}}}

\title{TeleCon connection modes\\[0.3em]
\large What to select before you connect}
\author{Phone setting, board, and sketch must match}
\date{Version 1.0 \quad 2026}

\begin{document}
\maketitle

\begin{abstract}
A connection mode is the link the phone opens to the ESP32. Pick it in the
application settings before you tap Connect. The sketch flashed on the board
must use the same handshake. This note lists every mode the app still offers,
and the camera overlay, which is not a separate mode.
\end{abstract}

{\sffamily\scriptsize
\chip{TblKnob}{Camera overlay}}

\tableofcontents
\vspace{1em}

\section{Set the mode before you connect}

The phone does not change the sketch when you change the mode. It only
changes how it speaks. Choose the mode, then connect. Classic Simple,
Classic Binary, and BLE Binary are all available.

\begin{itemize}[leftmargin=1.4cm]
  \item The mode stays in effect for that application until you change it.
        You do not reconnect just to remap a telemetry channel. You do
        reconnect, and you reflash, when the handshake itself changes.
\end{itemize}

\section{Modes at a glance}

\begin{center}
\small
\renewcommand{\arraystretch}{1.25}
\setlength{\tabcolsep}{4pt}
\begin{tabular}{>{\raggedright\arraybackslash}p{4.2cm}
                >{\raggedright\arraybackslash}p{2.4cm}
                >{\raggedright\arraybackslash}p{7.2cm}}
\rowcolor{TblHead}
\thc{Mode in the app} & \thc{Board} & \thc{Sketch} \\
\rowcolor{TblChan}
Classic + Simple & DevKit & \code{TeleCon_Classic_Simple} \\
\rowcolor{TblChan}
Classic + Binary & DevKit & \code{TeleCon_Classic_Binary} \\
\rowcolor{TblChan}
BLE + Binary & DevKit & \code{TeleCon_BLE_Binary} \\
\end{tabular}
\end{center}

The optional row \ui{Use SoftAP camera} sits on top of the three Bluetooth
modes. It does not replace them. See the last section.

\section{Classic + Simple}

Bluetooth Classic, text lines. This is the first mode to use.

\begin{itemize}[leftmargin=1.4cm]
  \item App: board \ui{DevKit}, mode \ui{Classic + Simple (text lines)}.
  \item Board name \code{ESP32-TC-RC-BT-Simple}, PIN \code{1234}.
  \item Handshake \code{RC:CONNECT,proto,simple}.
  \item Sticks, buttons, and telemetry are lines such as \code{RC:CTRL},
        \code{RC:BTN}, \code{RC:DATA}, and \code{RC:PLOT}.
  \item Eight telemetry bytes (CH1--CH8) when the sketch sends a plot.
\end{itemize}

\section{Classic + Binary}

Same Classic radio as Simple, with short binary frames instead of text.
Higher rate, same pins, less to read in the Serial Monitor.

\begin{itemize}[leftmargin=1.4cm]
  \item App: board \ui{DevKit}, mode \ui{Classic + Binary}.
  \item Board name \code{ESP32-TC-RC-BT-Binary}.
  \item Handshake \code{RC:CONNECT,proto,binary}.
  \item Frames start with \code{AA 55} (sticks), \code{BB 66} (buttons),
        and \code{CC 11/22/33} (telemetry).
  \item Do not pair this board as BLE, and do not leave the phone on
        Simple.
\end{itemize}

\section{BLE + Binary}

Bluetooth Low Energy, Nordic UART, with the same binary frames as Classic
Binary. Use it for a BLE-only module such as ESP32-C3 or ESP32-C6, or when
you want the lower idle radio of BLE.

\begin{itemize}[leftmargin=1.4cm]
  \item App: board \ui{DevKit}, mode \ui{BLE + Binary}.
  \item Board name \code{ESP32-TC-RC-BLE-Binary}.
  \item Handshake \code{RC:CONNECT,proto,binary}.
  \item There is no Classic PIN. Scan for the BLE name. Do not pair it as
        a Classic device.
\end{itemize}

\section{SoftAP camera is not a mode}

\ui{Use SoftAP camera} is a switch on a DevKit Bluetooth mode. The control
mode stays Classic Simple, Classic Binary, or BLE Binary.

\begin{itemize}[leftmargin=1.4cm]
  \item Flash \code{TeleCon_CAM_SoftAP_Video} on a separate ESP32-CAM.
  \item Join \code{TeleCon-RC-CAM-Starter} / \code{telecon1234} for the
        picture only.
  \item Sticks stay on the Bluetooth mode you already selected.
  \item The phone does not open TCP port 3333 for this overlay.
\end{itemize}

One ESP32-CAM does not carry the sticks. A DevKit does Bluetooth, and a
second CAM does video only.

\end{document}
